\section{Cognitive Map：空间只是结构推理的一个特例}

认知地图最初来自 navigation，但课程的目标是把它扩展到任意 relational structure。我们先看空间证据，再看非空间 transitive inference，最后把共同问题抽象为 graph learning：观察内容变化，边关系保持，模型需要重用结构并快速更新内容绑定。

\subsection{从 Tolman 的地图到树状泛化}

Theory/navigation 页把 cognitive map 描述为对环境结构的内部模型，使 agent 能在道路阻断或目标改变时重新规划，而不是只复现一条 cached action sequence。树状例子进一步说明，结构不一定是二维平面；层级、亲属关系和任务状态都可以形成 graph，关键是模型能否推断未直接体验的路径。

\lecturefigure{slide-43-cognitive-map-theory-navigation.jpg}{Cognitive-map 理论：结构模型支持灵活导航}{官方录像手写教学状态 V043；00:47:32}

\lecturefigure{slide-44-tree-structured-generalization.jpg}{树状关系中的结构泛化}{官方录像手写教学状态 V044；00:48:20}

若环境表示为图 $G=(V,E)$，每个节点 $i$ 有观察 $x_i$，边类型/动作 $a_t$ 决定 transition $i_t\to i_{t+1}$。结构泛化要求学习
$$
g_{t+1}=f_{\theta}(g_t,a_t),
$$
其中 $g_t$ 是与具体 observation identity 尽量解耦的 structural state。换一组 $x_i$ 时，$f_{\theta}$ 仍可更新位置/关系；只需重新绑定 $g_i$ 与新内容。

\begin{knowledgebox}{背景概念：cognitive map 不等于地理地图}
地图的本质是关系可组合：A 在 B 左边、B 在 C 左边，可推出 A 与 C 的相对关系；某动作从节点 1 到节点 2，同一动作规则可迁移到新对象。二维空间是最直观案例，但 graph、family tree、任务状态和概念关系都可使用同一结构思想。
\end{knowledgebox}

\subsection{Place cells、grid cells 与空间结构证据}

Evidence overview 页汇集空间任务、neural recording 和行为结果。讲者口头回顾 place cell：某个 hippocampal neuron 在动物处于环境特定位置时强烈 firing；grid cell 则在多个位置形成 hexagonal lattice。它们不是直接的 GPS 坐标，而是神经活动随位置呈现稳定结构。

\lecturefigure{slide-45-task-evidence-selection.jpg}{认知地图的行为与神经证据概览}{官方录像手写教学状态 V045；00:51:44}

\teachervoice{Will 用 London taxi drivers 的“Knowledge”考试说明空间结构学习的行为代价，并引用 hippocampal size 与驾驶经验的相关证据。他随即转向非空间任务，因为“空间相关”还不能证明脑区表示的是更一般的 relational structure。}

\begin{warningbox}{神经 tuning curve 的常见误读}
一个 cell 在特定位置 firing，不等于它单独“存储这个位置”；位置通常由 population code 表示。grid pattern 与模型 position encoding 相似，也不代表训练过程、输入来源和生物功能完全相同。应比较可解码信息、变换规律与任务因果作用。
\end{warningbox}

\subsection{非空间证据：Transitive inference}

Transitive-inference 页展示一类非空间关系任务：动物学习 A 比 B、B 比 C 等局部 pair，再被测试未直接训练的 A 与 C。若只记忆 observed pairs，无法稳定推断新组合；若形成 ordered relational map，就能沿结构推导。手写页把行为实验、成像证据和 pair relation 放在同一条论证链上。

\lecturefigure{slide-46-transitive-inference-evidence.jpg}{Transitive inference：从已学 pair 推断未见关系}{官方录像手写教学状态 V046；00:54:54}

\lecturefigure{slide-47-executive-layer-map.jpg}{关系任务中的 observation、structure 与行为读出}{官方录像手写教学状态 V047；00:55:18}

\begin{knowledgebox}{读图：三种策略怎样区分}
\term{pair memorization} 只能回答训练过的 pair；\term{scalar ordering} 给每个对象一个一维 rank，适合线性序；\term{relational graph} 表示多种 edge，能处理树、环与二维空间。实验设计要用未见组合和结构改变区分这些策略，而不是只测训练 pair accuracy。
\end{knowledgebox}

\subsection{从 graph world model 到通用神经算法}

Graph-world 页把环境抽象为节点、边和 observation；下一页列出目标算法：结构应可跨环境复用，内容要快速写入，规划应在 latent graph 上进行，神经表征还应与 hippocampal/entorhinal data 对照。这里的挑战同时包含 representation learning、memory binding 和 model-based inference。

\lecturefigure{slide-48-graph-world-model.jpg}{把环境抽象为可复用 graph 与可变化 observations}{官方录像手写教学状态 V048；00:58:58}

\lecturefigure{slide-49-neural-algorithm-general.jpg}{理想神经算法：结构迁移、快速绑定与推理}{官方录像手写教学状态 V049；00:59:28}

\begin{importantbox}{第二讲的任务分解}
慢学习负责掌握 transition structure 与可组合规则；快学习负责把新 sensory content 绑定到 structural states；推理负责沿结构传播与规划。TEM 的架构正是把这三项分给不同表示和更新机制。
\end{importantbox}

\subsection{本章小结}

Place/grid cells 和 taxi-driver evidence 支持空间认知地图，transitive inference 等任务则支持更广 relational map。共同计算问题是：把 observation identity 与 reusable structure 分开，再通过快速记忆把二者绑定，使未见组合也可推断。

